\documentclass[a4paper,12pt]{article}
\usepackage[utf8]{inputenc}
\usepackage[T2A]{fontenc}
\usepackage[russian]{babel}
\usepackage[margin=2cm]{geometry}
\usepackage{parskip}

\title{Текст выступления к защите ВКР}
\author{Горбачев М.\,Д., М01-404}
\date{Май 2026}

\begin{document}
\maketitle

\textit{Ориентир: 7 минут. Метки }\texttt{[Слайд N]}\textit{ — момент перелистывания.}

\section*{[Слайд 1 — Титул]}

Здравствуйте, уважаемая комиссия. Меня зовут Горбачев Максим, группа М01-404. Тема моей работы — поиск одного и того же автомобиля на видео с разных камер наблюдения. Научный руководитель — Назаров Алексей Николаевич.

\section*{[Слайд 2 — Постановка задачи]}

Представьте: есть фотография машины с одной камеры. Нужно найти эту же машину на снимках с других камер города. Номер мы не используем — он может быть грязным, нечитаемым, или его съёмка запрещена законом. Программа должна узнавать машину только по внешнему виду. Важная деталь: машины на этапе обучения и на этапе проверки — разные. То есть программа должна научиться сравнивать автомобили в принципе, а не запоминать конкретные.

\section*{[Слайд 3 — Сложности]}

Почему это трудно. Разные камеры снимают под разными углами и с разного расстояния. Освещение меняется в течение дня. Машину может частично закрывать другая машина или столб. И самое главное — автомобили выпускаются серийно: тысячи одинаковых седанов одного цвета. Уникальных черт, как у человека, у них нет.

\section*{[Слайд 4 — Датасет]}

Для обучения я взял открытый набор данных VeRi-776. Около 50 тысяч фотографий, 776 автомобилей, снятых на 20 камер. Справа на слайде — схема, какая камера с какой связана дорогами.

\section*{[Слайд 5 — Топология]}

Камеры образуют реальную городскую сеть. В среднем одна машина попадает в кадр девяти камер. И сразу важное наблюдение: одна из камер, под номером 18, даёт почти нулевой результат у всех моделей. Это не ошибка программы — это плохое качество самой камеры. К этому вернёмся.

\section*{[Слайд 6 — Первая модель]}

Я сравнивал две группы моделей. Первая — классическая, она называется ResNet. Это свёрточная нейросеть, она смотрит на картинку постепенно, начиная с мелких деталей и собирая их в более крупные. Брал две её версии: поменьше и побольше. Большая версия — на случай если объёма хватит для более сложной модели.

\section*{[Слайд 7 — Вторая модель]}

Вторая модель — более современная, она называется Vision Transformer. Она устроена иначе: режет картинку на маленькие квадраты и сразу смотрит, как все квадраты связаны друг с другом. Это даёт ей широкий взгляд на изображение целиком. И что важно — эта модель была предобучена на куда большем наборе картинок: 14 миллионов против одного.

\section*{[Слайд 8 — Как обучается]}

Обучение строится на простом принципе: показываем модели три картинки сразу. Одна и та же машина с двух камер и любая другая машина. Модель должна научиться: первые две — близкие, третья — далёкая. Плюс параллельно она учится угадывать, какая именно это машина из обучающего набора.

\section*{[Слайд 9 — Условия обучения]}

Обучал в облаке Google Colab на бесплатной видеокарте. На таблице — настройки. Главное: обе модели обучались в одинаковых условиях, чтобы сравнение было честным.

\section*{[Слайд 10 — Результаты]}

Главный слайд работы. Метрика \textit{Rank-1} — это доля случаев, когда самая первая выданная машина совпала с правильной. \textit{mAP} — общее качество поиска по всему списку. Vision Transformer выигрывает: правильный ответ первым в 90\% случаев. Большая ResNet оказалась хуже маленькой — это типичный случай, когда модели слишком много, а данных слишком мало, и она просто запоминает обучение. Дополнительно применил приём, который называется re-ranking — он пересортировывает результаты с учётом соседей, и качество ещё подросло.

\section*{[Слайд 11 — Картинки эмбеддингов]}

Чтобы понять, почему Vision Transformer лучше, я посмотрел, как модели располагают машины в своём внутреннем пространстве. Каждая машина превращается в набор чисел; чем числа похожее — тем машины похожее. У ResNet точки одного автомобиля разбросаны и перемешаны с другими. У Vision Transformer — собраны в плотные группы.

\section*{[Слайд 12 — Числа про эмбеддинги]}

В числах это видно ещё нагляднее. Машины одного типа у Transformer лежат почти в 2 раза кучнее. Разные машины — в полтора раза дальше друг от друга. А общая оценка качества разделения у Transformer лучше в десять раз.

\section*{[Слайд 13 — На что смотрит модель]}

Дальше я хотел понять, на что именно смотрит Transformer. У него есть встроенный механизм внимания — можно вытащить и нарисовать. Простой способ дал размытую картинку, поэтому я взял внимание из последних слоёв и усреднил по-другому — результат стал чище.

\section*{[Слайд 14 — Карта внимания]}

Вот результат. Яркие области — то, на что модель обращает внимание. Видно: она смотрит на кузов и характерные детали, а фон игнорирует. То есть она действительно учится распознавать машины, а не дороги.

\section*{[Слайд 15 — Удачный пример]}

Пример удачной работы. Запрос — жёлтый хэтчбек. Обе модели находят все пять правильных снимков из топ-5.

\section*{[Слайд 16 — Неудача обеих]}

А вот пример, где ни одна модель не справилась. Чёрный внедорожник в темноте. Все пять выдач — другие машины. Это предел, который не зависит от модели: на тёмном снимке низкого качества не видно никаких характерных черт.

\section*{[Слайд 17 — Где Transformer лучше]}

А вот случай, где разница видна. Оранжевый хэтчбек. ResNet нашёл одну машину из пяти. Transformer — четыре из пяти. Здесь ему помог широкий взгляд: он одновременно учитывает форму, цвет и фактуру кузова.

\section*{[Слайд 18 — Полная система]}

Чтобы показать применимость, я собрал готовую систему. Камера снимает кадр. Быстрый детектор находит на нём машины. Каждая машина отправляется в Transformer, превращается в набор чисел. Эти числа сравниваются с базой через быстрый поиск. Полная обработка кадра — доли секунды. На реальной нарезке детектора качество чуть падает: 86 процентов против 90 на идеальных снимках. Это цена работы с настоящим видео.

\section*{[Слайд 19 — Покамерная карта]}

Здесь видно, как модель работает между разными парами камер. Зелёное — хорошо, красное — плохо. Видно ту самую камеру 18 — сплошной красный столбец. Если бы я смотрел только средний результат, эту проблему я бы не заметил. Это важный практический вывод: общая цифра скрывает локальные сбои.

\section*{[Слайд 20 — Отрицательный результат]}

Один эксперимент не оправдал ожиданий, но это тоже результат. Я попробовал более «умный» способ сравнивать наборы чисел, который учитывает их взаимосвязи. На одной модели стало чуть лучше, на Transformer — резко хуже, минус 31 процент. Разобрался почему: после нормировки числа лежат на поверхности сферы, а этот способ сравнения работает с другой геометрией. Они несовместимы. Вывод: для нормированных наборов чисел простое сравнение по углу — наилучший выбор.

\section*{[Слайд 21 — Что сделано]}

Что в итоге сделано: обучены три модели, реализовано переупорядочивание результатов, проведён анализ внутреннего устройства моделей, разобран случай неисправной камеры, собрана работающая система обнаружения и поиска.

\section*{[Слайд 22 — Выводы]}

Главные выводы. Современная модель Vision Transformer на этой задаче работает лучше классической свёрточной сети — за счёт более широкого взгляда и большего объёма предварительного обучения. Просто увеличить размер модели не помогает — нужно больше данных. Средние цифры могут скрывать локальные проблемы — нужно смотреть в разрезе. И, наконец, не любая «умная» математика улучшает результат — иногда простое решение оптимально по самой природе задачи.

Из перспектив: пробовать ещё более крупные современные модели, добавить отслеживание машин во времени и адаптировать систему под российские камеры.

\section*{[Слайд 23 — Спасибо]}

Спасибо за внимание. Готов ответить на вопросы.

\end{document}
